% Pista alg-25s-q2 · Àlgebra
% Procedència: PAU Catalunya 2025 · Convocatòria de setembre · Sèrie 3 · Exercici 2
\documentclass[11pt,a4paper]{article}
\usepackage{pista}
\pistainfo{Catalunya 2025}{Convocatòria de setembre}{Sèrie 3}

\begin{document}

\section*{\textcolor{blauPista}{Exercici 2}}

\noindent Considereu el sistema d'equacions lineals següent:
\[
\left.\begin{array}{r}
x + 3y + z = 5 \\
mx + 2z = 0 \\
my - z = m
\end{array}\right\}
\]

\subsection*{\textit{a)} \normalfont Discutiu el sistema per als diferents valors del paràmetre $m$. \hfill\textnormal{[1,25~punts]}}

\pista{Escriu la matriu de coeficients $A$ i la matriu ampliada $A^{'}$. Calcula $\det(A)$ i troba els valors del paràmetre $m$ que l'anul·len. Per a cadascun d'aquests valors, compara els rangs d'$A$ i d'$A^{'}$. No t'oblidis del cas en què $\det(A) \neq 0$: quin tipus de sistema és?}

\subsection*{\textit{b)} \normalfont Resoleu el sistema per a $m = 1$. \hfill\textnormal{[0,5~punts]}}

\pista{Substitueix $m = 1$. Segons la discussió de l'apartat a), quin tipus de sistema obtens? Per resoldre'l, aïlla $x$ de la segona equació i $y$ de la tercera en funció de $z$, i substitueix-les a la primera.}

\subsection*{\textit{c)} \normalfont Resoleu el sistema quan aquest tingui infinites solucions. \hfill\textnormal{[0,75~punts]}}

\pista{De l'apartat a) ja saps quin valor del paràmetre fa que el sistema sigui compatible indeterminat. Fixa aquest valor i resol el sistema introduint una variable lliure (per exemple $y = \lambda$) per expressar la resta de variables en funció d'ella. Compte: amb aquest valor, les dues últimes equacions ja determinen $z$, de manera que $z$ no pot ser la variable lliure.}

\end{document}
